\chapter{Wstęp}
\label{chap:introduction}
\section{Kontekst i motywacja}
Projekt EdgeGuard IoT stanowi planowane rozwinięcie wcześniejszego prototypu inteligentnego domu. Nowa wersja zakłada rozproszenie czujników i elementów wykonawczych pomiędzy niezależne węzły ESP32 oraz analizę telemetrii na bramie Raspberry Pi.
\writingnote{Dodaj odwołanie do własnej pracy inżynierskiej i opisz jej zakres. Wyraźnie wskaż ponownie wykorzystane elementy oraz nowy wkład magisterski.}

\section{Cel i problem badawczy}
Celem pracy jest zaprojektowanie systemu oraz zbadanie przydatności lekkich metod wykrywania anomalii w środowisku inteligentnego domu przy ograniczonych zasobach bramy brzegowej.
\writingnote{Uzgodnij z promotorem dokładne brzmienie celu i zakres wykrywanych nieprawidłowości.}

\section{Pytania badawcze}
\begingroup\small\singlespacing
\begin{enumerate}[label=RQ\arabic*:,leftmargin=1.5cm]
\item Która metoda zapewnia najlepszy kompromis między skutecznością detekcji a zapotrzebowaniem na zasoby Raspberry Pi?
\item Czy modele trenowane głównie na zgodnych danych publicznych lub syntetycznych wykrywają porównywalne anomalie w danych fizycznego prototypu?
\item Jak skutecznie można wykrywać anomalie związane z poszczególnymi węzłami?
\item Jak liczba węzłów fizycznych lub symulowanych wpływa na skuteczność i zużycie zasobów?
\item W jakich warunkach ML daje korzyść względem prostych reguł?
\end{enumerate}
\endgroup

\section{Zakres i organizacja pracy}
\writingnote{Opisz zakres obowiązkowy, rozszerzenia opcjonalne oraz zawartość kolejnych rozdziałów. Oddziel przygotowanie infrastruktury od eksperymentalnej oceny metod.}
